\section{Statement and relation to known work}

For $r>0$ let $\Sph(r)$ carry the graph whose vertices are the points of the
sphere of radius $r$ and whose edges join points at Euclidean distance $1$
(chords, not geodesics). Write $h$ for the largest number of points on $\Sph(1)$
whose chord-$1$ graph has independence number at most $2$. Following Bachoc, Bekker, Moustrou and de Oliveira Filho~\cite{Bachoc2021}, a set of points on the unit sphere in
$\R^n$ is \emph{$t$-almost equiangular} if among any three of them some pair has
inner product $t$, and $\indep(n,t)$ denotes the maximum cardinality of such a
set. A chord of length $1$ on the unit sphere means inner product $\tfrac12$, so
\[
h \;=\; \indep(3,\tfrac12).
\]

\begin{theorem}
\label{thm:h7}
$h=7$, with an extremal configuration unique up to isomorphism ($11$ edges,
degree sequence $[3,3,3,3,3,3,4]$) and realisable over $\Q(\sqrt{46})$.
\end{theorem}

Two neighbouring lines of work stop short of this value. Bachoc et al.~\cite{Bachoc2021} give a
semidefinite bound for $\indep(n,t)$ and prove $\indep(n,t)\le 2(n+1)$ for
$t\le0$; our $t=\tfrac12$ is positive and outside that range. Rosenfeld~\cite{Rosenfeld1991} showed
that an almost-equidistant set on a sphere of radius $1/\sqrt2$ in $\R^d$ has at
most $2d$ points, and Bezdek and L\'angi~\cite{BezdekLangi} extended this to $r<1/\sqrt2$ with the
bound $2d+2$; our radius is $r=1>1/\sqrt2$, again outside the range covered. The
literature recorded only $3\le h\le8$.

\begin{remark}[what is and is not new, stated plainly]
The tables of the semidefinite bounds were checked~\cite{Bachoc2025}. Bachoc et al.~\cite{Bachoc2021} treat
$t\in[-1,0]$ throughout, so $\indep(3,\tfrac12)$ is not among their computed
values, and they observe that the bound $2(n+1)$ fails for positive $t$.
However their Theorem~5.2~\cite{Bachoc2021} does settle the lower bound: $\mathrm{MS}(2)$ is
$(3,t)$-realizable for $t\ge t_{2,3}=0.2553\ldots$, so $h\ge7$ follows from
their work, and the extremal configuration below is the Moser spindle. The
contribution of this paper is the upper bound $h\le7$, the uniqueness, and the
exact Gram realisation.
\end{remark}

\paragraph{Wider context.} Chromatic numbers of spheres in their own right are
surveyed by Malen~\cite{Malen2014}. Repeated and unit distances on spheres go
back to Erd\H{o}s, Hickerson and Pach's work on Leo Moser's problem
~\cite{ErdosHickersonPach1989}, with further bounds by Swanepoel and
Valtr~\cite{SwanepoelValtr}; Soifer's survey~\cite{Soifer2008} places these
questions within the broader Hadwiger--Nelson literature. The cone route
discussed below is the one used to obtain the current best space bound
$\chrom(\R^3)\ge6$, via de Grey's $59$-vertex unit-distance
graph~\cite{deGrey2020R3}, itself building on the Polymath16
project~\cite{Polymath16}.

\section{Why the sphere}

One route to $\chrom(\R^3)\ge7$ runs through the \emph{cone}: if $A$ is a point
adjacent to everything else, then $\chrom(\R^3)\ge 1+\chrom(\Sph(1))$, since the
neighbours of $A$ lie exactly on the unit sphere about it. Section~\ref{sec:niven}
shows the cone route is closed, but the spherical questions it raises are not.

The quantity matters for a second reason. Nechushtan's chain in $\R^3$~\cite{Nechushtan2002} hangs on
the equality $n=5\indep$ for a $10$-point set, and that $10$ is the maximum size
of an almost-equidistant set in $\R^3$~\cite{Gyorey2020} --- a proved extremal fact. The spherical
analogue of that extremal fact is exactly $h$, and it feeds directly into the
upper bound of the bound of the companion report.

\section{The upper bound $h\le8$}

\begin{lemma}[$K_4$-freeness]
\label{lem:k4free}
The chord-$1$ graph on $\Sph(1)$ contains no $K_4$.
\end{lemma}

\begin{proof}
A chord of length $1$ on the unit sphere subtends $60^\circ$, so adjacency means
inner product $\tfrac12$. Four mutually adjacent unit vectors would have Gram
matrix $\tfrac12(I+J)$, which is positive definite, hence they would be linearly
independent --- impossible in $\R^3$.
\end{proof}

\begin{lemma}
For a graph $G$, $\indep(G)\le2$ if and only if the complement of $G$ is
triangle-free.
\end{lemma}

\begin{theorem}
$h\le8$.
\end{theorem}

\begin{proof}
A set realising $h$ gives a graph with no $K_4$ and no independent set of size
$3$. Since $R(3,4)=9$, such a graph has at most $8$ vertices.
\end{proof}

The recorded bound $h\le8$ was previously an assumption; the argument above
makes it a proof, and it does more: the $(3,4)$-Ramsey graphs on $8$ vertices
are \emph{unique} up to isomorphism --- the complement of the Wagner
(M\"obius--Kantor) graph $V_8$. So at $n=8$ there is exactly one combinatorial
candidate to rule out.

\section{Exclusion of $n=8$}

Let $G$ be a candidate and $\{a,b,c\}$ a triangle of $G$. The three
corresponding unit vectors have Gram matrix $M=\tfrac12(I+J)$, which is
invertible, so they form a basis of $\R^3$. For any other vertex $v$ put
$t_v=(\langle v,a\rangle,\langle v,b\rangle,\langle v,c\rangle)^{\mathsf T}$.
Then for all $u,v$,
\[
\langle u,v\rangle=t_u^{\mathsf T}M^{-1}t_v,\qquad
|v|^2=t_v^{\mathsf T}M^{-1}t_v .
\]
The entries of $t_v$ equal $\tfrac12$ where $v$ is adjacent to a triangle vertex
and are unknowns otherwise. Realisability becomes a polynomial system: one
equation $|v|^2=1$ per remaining vertex, one equation $\langle
u,v\rangle=\tfrac12$ per remaining edge. \emph{Non-edges are left entirely
free} --- no inequality is imposed on them --- so an empty solution set is
genuine non-realisability rather than an artifact of over-constraining.

\begin{theorem}
The complement of the Wagner graph is not realisable as a chord-$1$ graph on
$\Sph(1)$. Consequently $h\le7$.
\end{theorem}

\begin{proof}
The graph is $4$-regular with $16$ edges. Choosing the triangle $\{0,2,5\}$, the
elimination above gives $12$ equations in $9$ unknowns over $\Q$. Its Gr\"obner
basis (grevlex) is $\{1\}$, so the system has no solution even over $\CC$; a
fortiori none over $\R$.
\end{proof}

Note the shape of the argument: the Ramsey bound did not merely restrict the
size, it produced a \emph{single} graph, and uniqueness is what makes one
Gr\"obner computation sufficient.

\section{The construction at $n=7$}

Among the $13\,842$ labelled graphs on $7$ vertices with $\indep\le2$ and
$\omega\le3$ there are $9$ isomorphism classes, and exactly one is realisable.

\begin{theorem}
$h\ge7$. The extremal configuration is unique up to isomorphism.
\end{theorem}

\begin{proof}
Take the Gram matrix $G$ with $G_{ii}=1$, $G_{ij}=\tfrac12$ on the $11$ edges,
and on the ten non-edges the values
\[
-\tfrac13,\qquad \tfrac{37}{64},\qquad
\tfrac18\pm\tfrac{\sqrt{46}}{16},\qquad
\tfrac{23}{64}\pm\tfrac{3\sqrt{46}}{32}.
\]
Symbolic computation gives $\rank G=3$ exactly, with eigenvalues
$\tfrac{57}{32}$, $\tfrac{167}{64}\pm\tfrac{\sqrt{21881}}{192}$ and $0$ with
multiplicity $4$ --- all non-negative, so $G$ is positive semidefinite and is
therefore the Gram matrix of seven vectors spanning $\R^3$. Every diagonal entry
is $1$, so the vectors are unit; the $11$ prescribed entries are $\tfrac12$, so
those pairs are at chord distance $1$; and no non-edge entry equals $\tfrac12$,
so the graph realised is exactly the intended one.
\end{proof}

\begin{remark}[on method]
The exact values were obtained by refining a least-squares realisation (residual
$2.2\times10^{-16}$) with Newton's method on the vanishing $4\times4$ minors
until $\max|\text{minor}|<10^{-56}$, then identifying each entry by integer
relation detection at $60$ digits. At $9$ digits the same detection returned a
\emph{spurious} quadratic with six-figure coefficients; algebraic identification
needs far more precision than the apparent accuracy of the input suggests.
Working with the Gram matrix rather than coordinates removes the rotational
gauge and reduces the unknowns from $21$ to $10$.
\end{remark}

\section{Two structural facts about the spherical family}

\subsection{The link, in closed form}
\label{sec:link}

A chord of length $1$ on a sphere of radius $r$ corresponds to angular distance
$\alpha$ with $2r\sin(\alpha/2)=1$, so $\cos\alpha = 1-\frac{1}{2r^2}$. The
neighbours of a vertex lie on a circle of angular radius $\alpha$; two of them
at azimuthal separation $\Delta$ are adjacent precisely when
\[
\cos\Delta \;=\; \frac{\cos\alpha}{1+\cos\alpha}\;=\;\frac{2r^2-1}{4r^2-1}.
\]

Three consequences are read straight off the formula. The link has
\emph{maximum degree $2$ at every radius}. The formula is defined exactly when
$r\ge1/\sqrt3$; at $r=1/\sqrt3$ one has $\cos\Delta=-1$ (the link is a perfect
matching), and as $r\to\infty$, $\Delta\to60^\circ$ and the link becomes a
hexagon --- the plane as a limit. Hence $\Delta$ runs \emph{monotonically}
through $[60^\circ,180^\circ]$ as $r$ runs through $[1/\sqrt3,\infty)$.

\begin{proposition}[which radii close the link]
\label{prop:radii}
The link contains a cycle exactly when $\Delta/2\pi\in\Q$, in which case it is a
disjoint union of $k$-cycles with $\Delta=2\pi j/k$, $\gcd(j,k)=1$. By the
monotonicity above the condition is $j/k\in[\tfrac16,\tfrac12]$. Every rational
in that interval gives a closing radius; there are countably many, dense in
$[1/\sqrt3,\infty)$.
\end{proposition}

\begin{center}
\begin{tabular}{@{}llll@{}}
\toprule
$j/k$ & $\Delta$ & $r$ & comment \\
\midrule
$1/2$ & $180^\circ$ & $0.5773503$ & link is a matching \\
$1/3$ & $120^\circ$ & $0.6123724$ & tetrahedron \\
$1/4$ & $90^\circ$ & $0.7071068$ & octahedron \\
$1/5$ & $72^\circ$ & $0.9510565$ & icosahedron \\
$2/5$ & $144^\circ$ & $0.5877853$ & link $C_5$, pentagram step \\
$\mathbf{2/7}$ & $102.857^\circ$ & $\mathbf{0.6503886}$ & \textbf{link $C_7$} \\
$\mathbf{2/11}$ & $65.455^\circ$ & $\mathbf{1.3144594}$ & link $C_{11}$,
radius $>1$ \\
$1/6$ & $60^\circ$ & $\infty$ & the plane, giving $H$ \\
\bottomrule
\end{tabular}
\end{center}

\begin{remark}[a correction to an earlier draft]
An earlier version of this work enumerated only the family $j=1$ and concluded
that ``exactly three spheres close, the circumspheres of the three regular
deltahedra''. That holds only for step $1$. With step $j\ge2$ the link is a star
polygon, but \emph{as an abstract graph} it is still $C_k$, and only that
matters for a chromatic argument. There are infinitely many closing radii, some
with $r>1$.
\end{remark}

\subsection{At $r=1$ the link is never a cycle}
\label{sec:niven}

\begin{theorem}[Niven~\cite{Niven1956}]
\label{thm:niven}
On $\Sph(1)$ the link of every vertex is a disjoint union of paths; it never
contains a cycle, at any scale.
\end{theorem}

\begin{proof}
At $r=1$ the formula of Section~\ref{sec:link} gives $\cos\Delta=(1/2)/(3/2)=1/3$. By
Niven's theorem~\cite{Niven1956}, if $\theta/\pi$ is rational and $\cos\theta$ is rational then
$\cos\theta\in\{0,\pm\tfrac12,\pm1\}$. Since $1/3$ is not among these,
$\Delta/\pi$ is irrational and $k\Delta\ne2\pi m$ for all integers $k,m$. Since
the link has maximum degree $2$, a cycle would require exactly such a return.
\end{proof}

This closes the cone route: the cone requires radius exactly $1$, while the good
wheel gadgets live at other radii. The two requirements are incompatible, and
the incompatibility is proved rather than observed. It also explains,
retrospectively, why densifying spherical graphs at $r=1$ left $\chrom$ stuck at
$4$ in every experiment: there was no primitive gadget to begin with.

Sharpened statement: $r=1$ is not merely ``a radius without a gadget''; it is
the radius where $\cos\Delta=1/3$, a rational value outside Niven's set. The
obstruction is \emph{number-theoretic}, not geometric.

\section{Open questions}

\begin{problem}
$\chrom(\Sph(1))$. We can certify $\chrom\ge4$ on explicit graphs of a few
hundred vertices, and Theorem~\ref{thm:niven} says the usual wheel-based machinery is
unavailable. Is $\chrom(\Sph(1))=4$?
\end{problem}

\begin{problem}[higher dimensions]
The argument of Lemma~\ref{lem:k4free} generalises verbatim: on $S^{d-1}(1)\subset
\R^d$ the chord-$1$ graph is $K_{d+1}$-free, so
$h_d\le R(3,d+1)-1$, giving $8,13,17,22$ for $d=3,4,5,6$. At $d=3$ the bound is
off by exactly one. Is it off by one in general, and is the extremal
configuration always unique? The elimination generalises too --- a $d$-clique is
a basis of $\R^d$ --- but the number of Ramsey graphs to rule out is no longer
$1$, and that uniqueness is what made the $d=3$ computation a single Gr\"obner
basis.
\end{problem}

\begin{problem}
Uniqueness at $n=7$ is verified here by exhaustion over the nine isomorphism
classes. Is there a direct argument?
\end{problem}

\section{Appendix: ancillary scripts}
\label{sec:ancillaryB}

\begin{center}
\begin{tabular}{@{}lp{7.6cm}@{}}
\toprule
script & what it asserts \\
\midrule
\texttt{h\_sphere.py} & the $n=8$ candidate is the complement of the Wagner
graph; $\indep\le2$, $\omega\le3$ verified \\
\texttt{h7.py} & $13\,842$ labelled graphs at $n=7$, $9$ isomorphism classes,
exactly one realisable \\
\texttt{h7\_gram.py} & Newton refinement to $\max|\text{minor}|<10^{-56}$;
integer-relation identification at $60$ digits \\
\texttt{h7\_theorem.py} & exact Gram over $\Q(\sqrt{46})$: rank $3$, positive
semidefinite, no non-edge at $\tfrac12$ \\
\texttt{h8\_impossible.py} & the triangle-basis elimination and the Gr\"obner
basis $\{1\}$ \\
\bottomrule
\end{tabular}
\end{center}

All computations are exact except where explicitly stated; floating point is
used only to propose candidates, never to decide one.

%%%%%%%%%%%%%%%%%%%%%%%%%%%%%%%%%%%%%%%%%%%%%%%%%%%%%%%%%%%%%%%%%%%%%%%%%%%%%%
